\section{Groupoid Representation of Stack}\label{sec:algebraicgeometry-groupoid}

   Recall that given a presheaf of groupoid $F:C^{op}\to \Grpd$, the Grothendieck construction of the presheaf determines the category fibered in groupoids associated to $F$.
   Let $F_0$ denote the object presheaf and $F_1$ the morphism presheaf,
   \begin{proposition}
    There is pullback diagram of category of fibered groupoids
    \begin{center}
        \begin{tikzcd}
F_1 \arrow[d] \arrow[r] & F_0 \arrow[d] \\
F_0 \arrow[r]           & F            
\end{tikzcd}
    \end{center}
   \end{proposition}
   \begin{proof}
        This can be done in the presheaf of groupoids, which simplifies down to the $2$-pullback pointwise at each object $c\in C$.
   \end{proof}
   \begin{corollary}
    The property descends to stacks, if we denote $F_i$ to be the stackification of the corresponding category fibered in groupoids.
   \end{corollary}

   Let $(c_0, c_1)$ be a groupoid object in $C$, then $(c_0,c_1)$ determines a category fibered in groupoids determined by $(\Hom(-,c_0), \Hom(-,c_1))$, hence a stack. I like to denote it as $[c_0/c_1]$. Suppose we have a stack $\M$, then a \textbf{groupoid representation} of $\M$ is a groupoid object in $C$ and equivalence of stack
   \[
   [c_0/c_1]\simeq \M
   \]
   We refer to $c_\bullet:=c_1\times_{c_0}c_1\times_{c_0}...\times_{c_0}c_1$ to be the \textbf{Cêch nerve} of $\M$.

   If the topology of $C$ is subcanonical, then the stackification process is unnecessary.

   \begin{lemma}
    Let $c_0 \in C$, $\M$ a stack over $C$, and there is representable morphism $c_0\to \M$, denote $c_1:=c_0\times_{\M}c_0$ the object representing the pullback. Then there is a groupoid structure on $(c_0,c_1)$. Moreover, if $c_0\to \M$ is essentially surjective that is a covering, that is for every $x\in C$ and $u\to \M$, the pullback $x\times_{\M}c_0 \to x$ determines a cover in $C$, then there is equivalence
    \[
    [c_0/c_1]\simeq \M
    \]
   \end{lemma}
   \begin{proof}
    We recall that given a morphism $a\to b$, there is a simplicial object $a_\bullet$ determined by $a_n:=a\times_ba\times_b ...\times_b$. The stack properties implies the groupoid property.

    The argument for the representation is quite abstract and in my opinion technical and not too important for our concern, which is chromatic homotopy. I refer to \cite[Lemma 8.4.8, Lemma 94.16.1]{AGStacksProject}, which latter even though assume specifically the stack over fppf-topology, but works in more general setting.
   \end{proof}
   
   Suppose $\M$ is a fpqc-stack over a scheme. Let $U$ be a $S$-scheme with $U\to \M$ a representable faithfully flat morphism. Then the essential surjectiveness assumption is redundant (easy check that requires both representability and faithfulness).

   \begin{proposition}
    Let $\M$ be a fpqc-stack, with an affine (hense representable) faithfully flat morphism $\Spec(A)\to \M$ such that $\Spec(\Gamma)\cong \Spec(A)\times_{\M}\Spec(A)$, then there is equivalence
    \[
        \Qcoh_{\M}\simeq co\Module_{(A,\Gamma)}
    \]where $co\Module_{(A,\Gamma)}$ is the category of $(A,\Gamma)$-comodule.
   \end{proposition}
   \begin{proof}
    We use the cartesian module characterization of quasicoherent module.
    Notice the equivalence $\M\simeq [\Spec(A)/\Spec(\Gamma)]$. Let $\F$ be a quasicoherent module, denote $M:=\F(\Spec(A))$, the Cêch nerve 
    \[
    ...\to \Spec(\Gamma\otimes_A\Gamma)\to \Spec(\Gamma)\to \Spec(A)\to \M
    \]determines a $(A,\Gamma)$-comodule structure on $M$, since $\F(\Spec(\Gamma))\cong \Gamma\otimes_A M$.

    Therefore, there is a functor $ \Qcoh_{\M}\to co\Module_{(A,\Gamma)}$. 

    On the other hand, an object $x\in [\Spec(A)/ \Spec(\Gamma)]$ living over object $\Spec(B)$ corresponds to a morphism $\Spec(B)\to \Spec(A)$. Therefore, for a $(A,\Gamma)$-comodule $M$, we define
    \[
    \tilde{M}(\Spec(B)\to \Spec(A))=B\otimes_A M
    \]

    Check that these constructions are inverses of each other.
   \end{proof}

